\begin{initiativkarte}{HAWKS Racing}

  % Bild (optional)
  \IfFileExists{bilder/hawks_racing.jpg}{%
    \begin{center}
      \includegraphics[width=\linewidth,height=5cm,keepaspectratio]{bilder/hawks_racing.jpg}
    \end{center}
    \vspace{0.4cm}
  }{}

  % Kurzbeschreibung
  \textit{\textcolor{hawgray}{HAWKS Racing ist das Formula Student Team der HAW Hamburg. Wir entwickeln und bauen Rennfahrzeuge und sammeln dabei praxisnahe Erfahrung in Ingenieurwesen, Software und Organisation.}}

  \vspace{0.5cm}
  \hawrule

  % Beschreibung
  \textbf{\textcolor{hawblue}{Was machen wir?}}\\[0.3cm]

  HAWKS Racing e.V. ist das Formula Student Team der Hochschule für Angewandte Wissenschaften Hamburg. Wir entwickeln und fertigen Rennfahrzeuge, mit denen wir an internationalen Formula Student Wettbewerben teilnehmen.

  Dabei verbinden wir Theorie und Praxis:
  Von der Konstruktion über Elektronik und Software bis hin zur Fertigung und dem Rennbetrieb arbeiten Studierende gemeinsam an einem echten High-Performance-Fahrzeug.

  Der Verein bietet Studierenden die Möglichkeit, praktische Erfahrungen zu sammeln, die über das Studium hinausgehen. Dazu gehören Werkstattarbeit, Projektmanagement, Softwareentwicklung und die Zusammenarbeit in interdisziplinären Teams.

  Außerdem unterstützen wir unsere Mitglieder bei:
  \begin{itemize}
      \item praktischer Ingenieurserfahrung im Fahrzeugbau
      \item Einblick in reale Entwicklungsprozesse
      \item Bachelor- und Masterarbeiten im Projektkontext
      \item Vernetzung mit Industriepartnern
      \item persönlicher und fachlicher Weiterentwicklung
  \end{itemize}

  Die Formula Student ist ein internationaler Konstruktionswettbewerb, bei dem Teams aus aller Welt eigene Rennfahrzeuge entwickeln – mit Fokus auf Engineering, Business Plan und Performance.

  \vspace{0.5cm}

  % Tags
  \tag{Engineering}
  \tag{Automotive}
  \tag{Software}
  \tag{Teamwork}
  \tag{Wettbewerb}

  \vspace{0.6cm}
  \hawrule

  % Infozeilen
  \infozeile{\faUsers}{Zielgruppe: Studierende aus Informatik, Elektrotechnik, Maschinenbau und Design}
  \infozeile{\faGraduationCap}{Semester: Ab ca. 1.–2. Semester möglich}
  \infozeile{\faMapMarker*}{Ort: Fahrzeuglabor HAW Hamburg, Berliner Tor 9}
  \infozeile{\faEnvelope}{Kontakt: \href{mailto:info@hawksracing.de}{info@hawksracing.de}}
  \infozeile{\faGlobe}{Website: \href{https://hawksracing.de}{hawksracing.de}}

\end{initiativkarte}